% ============================================================
%  Chapter 6: The Propagation Bound and Lorentz Invariance
%  Source: bounded-phase-space-trajectory.tex, Consequences 16--18
% ============================================================

\epigraph{%
  The speed of light is not a postulate.\\
  It is the rate at which a bounded system\\
  can propagate information about its own partition structure.%
}{}

\section{The Propagation Bound}
\label{sec:ch6-propbound}

\begin{consequence}{Universal Propagation Bound $c$}{propagationbound}
  The rate at which the partition depth field $\Depth(\mathbf{x},t)$ can
  propagate through space is bounded above by a constant $c$:
  \begin{equation}
    \frac{|\Delta\mathbf{x}|}{\Delta t} \leq c,
  \end{equation}
  where $c = \bar{x}_\Omega / \bar{t}_\Omega$ is the ratio of the mean
  spatial extent of $\Omega$ to the mean recurrence time of the system.
\end{consequence}

\begin{proof}[Proof sketch]
  By Kac's lemma, the mean recurrence time of a set $A$ of measure $\mu(A)$
  under a measure-preserving flow on a space of total measure $\mu(\Omega)$ is
  $\bar{t}_A = \mu(\Omega)/\mu(A)$.
  The minimum recurrence time corresponds to the largest possible set
  $A = \Omega$ itself, giving $\bar{t}_\Omega = 1$ (in natural units for the
  system).
  The spatial extent of $\Omega$ is bounded by clause~(i).
  Their ratio defines an invariant speed $c$.
\end{proof}

\begin{remark}
  This is essentially the Ignatowsky derivation of the Lorentz group (1910)
  without the light postulate: from the finite extent of a bounded space and
  the Kac mean recurrence time, a finite invariant speed emerges.
  The identification of this speed with the speed of light in vacuum is
  an empirical identification, not a postulate.
\end{remark}

\section{Categorical Speed of Light}
\label{sec:ch6-catc}

\begin{consequence}{Categorical Speed}{categoricalc}
  In terms of partition coordinates, the invariant speed is
  \begin{equation}
    c = \frac{\delta x}{\taup},
  \end{equation}
  where $\delta x$ is the spatial resolution of one partition cell and $\taup$
  is the partition lag time (the minimum time to move between adjacent cells).
\end{consequence}

This gives $c$ a direct physical meaning: it is the maximum rate at which a
system can update its partition cell assignment.
An information signal cannot propagate faster than one partition cell per
partition lag time.

\section{Lorentz Invariance}
\label{sec:ch6-lorentz}

\begin{consequence}{Lorentz Invariance}{lorentz}
  The dynamics generated by the partition Lagrangian~\eqref{eq:partitionlagrangian}
  is invariant under the Lorentz group $SO(1,3)$.
\end{consequence}

\begin{proof}[Proof sketch]
  By the Ignatowsky group-composition argument:
  \begin{itemize}
    \item Spatial isotropy (the partition structure has no preferred direction
      in the bulk of $\Omega$) implies that velocity compositions form a group.
    \item The existence of a finite invariant speed $c$
      (Consequence~\ref{cons:propagationbound}) selects the Lorentz group
      uniquely from among the one-parameter family of groups compatible with
      isotropy and associativity.
    \item The discrete Koopman spectrum (Consequence~\ref{cons:discretemodes})
      supplies the boost parameter.
  \end{itemize}
  The Lorentz group is the symmetry of the partition dynamics.
\end{proof}

\begin{remark}
  Special relativity is typically founded on two postulates: the principle of
  relativity and the constancy of $c$.
  Here, both are consequences of a single axiom:
  \begin{itemize}
    \item The principle of relativity follows from the spatial isotropy of the
      bounded phase space (no preferred direction is encoded in BPSL).
    \item The constancy of $c$ follows from the Kac recurrence bound
      (Consequence~\ref{cons:propagationbound}).
  \end{itemize}
  Special relativity is therefore a corollary of BPSL.
\end{remark}

\section{The Dispersion Relation}
\label{sec:ch6-dispersion}

The Lorentz group and the discrete Koopman spectrum together yield the
dispersion relation.

\begin{consequence}{Relativistic Dispersion}{dispersion}
  Each Koopman mode with rest frequency $\omega_0$ satisfies
  \begin{equation}\label{eq:dispersion}
    \omega^2 = \omega_0^2 + c^2 k^2,
  \end{equation}
  where $k$ is the spatial wavenumber of the mode.
\end{consequence}

This is the relativistic energy--momentum relation
$E^2 = (m_0 c^2)^2 + (pc)^2$ in frequency--wavenumber language.
It has been derived without assuming the existence of particles, mass, or
the equivalence principle.

\bigskip
\begin{tcolorbox}[colback=crownblue!5, colframe=crownblue!60!black,
                  title={\textbf{Chapter 6 Summary}}]
\begin{itemize}[noitemsep]
  \item The Kac recurrence bound on a finite-measure phase space yields a
    finite invariant propagation speed $c = \delta x / \taup$.
  \item The Ignatowsky argument, applied to the partition dynamics, gives the
    Lorentz group as the unique symmetry compatible with isotropy and a finite
    invariant speed.
  \item Lorentz invariance of electrodynamics, special relativity, and the
    relativistic dispersion relation are all Consequences.
  \item The light postulate is not needed; the constancy of $c$ is derived.
\end{itemize}
\end{tcolorbox}
